\chapter{Reference-Informed Judge Generation}\label{app:judge-generation}

This appendix details the rubric-generation mechanics referenced in
Chapter~\ref{ch:experimental-design}.

\section{Generation Procedure}

The judge rubric generation procedure (Section~\ref{sec:arch-judge}) is
instantiated as follows. A leaf's induction pool is its construction
queries \emph{after} two filters: the reserved-set filter, which is audited
by a build-failing assertion and is what makes the answer-key-blind claim
enforceable rather than asserted, and a hash filter that thins the pool by
roughly a fifth.

That pool is used in full. \texttt{TaxonomyJudgeService.generateJudgeForNode}
chunks it into batches of $25$ (\texttt{details.chunked(25)}), induces partial
guidelines from \emph{every} batch, and synthesises them into one rubric. The
$25$ is a batch size set by prompt length, not a sample size: a leaf's rubric is
induced from all of its filtered construction queries, not from a representative
subset of them.
The map-step LLM and the reduce-step LLM are the same underlying model
as the arena judge, to avoid introducing a second
model's stylistic biases into the rubric. That choice has a cost, and
it is stated here because no other section carries it: inducing the
rubric with the model that will later apply it bakes the judge's own
priors into the criteria, which weakens the rubric-against-generic
contrast of Section~\ref{sec:res-c5}. A judge that agrees
with a rubric written by itself is a weaker test than a judge that
agrees with a rubric written by someone else. Guideline generation is
a one-time cost per frozen snapshot: rubrics are computed once at
snapshot-freeze time and reused for every subsequent arena run against
that snapshot, regardless of which models are being compared.

\section{Structural Template}\label{sec:judge-template-full}

The global structural template, covering presentation order, tie
handling, and output format, is fixed across all leaves. It governs
only the shared scaffolding of the judge prompt; the leaf-specific
reference-informed content (representative queries, answer-key
context) is injected into this template at generation time and does
not modify its structure.

The separation the template enforces is what makes the two arms
comparable, so it is stated explicitly. The induced, leaf-specific half
defines \emph{what} to evaluate; the fixed half defines \emph{how}. The
synthesis step that produces each leaf's judge
(\texttt{JudgePrompts.synthesizeFinalJudge}) is instructed in these
terms, and the instruction is quoted rather than paraphrased because
the constraint it imposes is the thesis's control:

\begin{quote}\small
``This persona defines WHAT to evaluate (domain knowledge, axioms,
failure modes) --- NOT HOW to evaluate. Verdict format, tie policy, and
evaluation steps are controlled externally and must NOT appear here.''
\end{quote}

\noindent It further requires each rubric to ``[s]pecialise strictly''
to its leaf and to ``ignore generic qualities that do not affect
technical correctness: formatting, markdown, response length, stylistic
elegance, model identity'', and forbids any ``compare A vs B'', ``pick
a winner'', or tie-break instruction from entering the induced text.
Verdict format is a JSON object with exactly two string keys,
\texttt{system\_prompt} and \texttt{rubric}.

The generic-rubric arm (C5) replaces the induced half with the MT-Bench
pairwise judge prompt verbatim (\texttt{GenericPairwiseJudgePrompt}).

\paragraph{The two arms differ in more than the rubric.} This is a limitation of
the contrast and is stated here because Chapter~\ref{ch:results} reports a
measured difference between the arms. The generic arm is not the rubric arm with
its system prompt swapped: it is a separate code path
(\texttt{TaxonomyArenaService}, the \texttt{isC5} branch), and it differs in four
respects at once.

\begin{table}[htbp]
\centering
\small
\caption{What actually differs between the two arms: four factors
travel together (system prompt, user message, output format,
confidence handling), while question set, stored responses, judge
model and protocol are identical by construction.}
\label{tab:arm-differences}
\begin{tabular}{p{2.6cm}p{4.6cm}p{4.6cm}}
\toprule
& Rubric arm (MAIN) & Generic arm (C5) \\
\midrule
System prompt & induced per leaf & MT-Bench, verbatim \\
User message & \texttt{[Question]} / \texttt{[Model A's Response]}, then a
mandatory four-step critique sequence & MT-Bench template, asking for ``a short
explanation'' \\
Output & structured JSON, schema-enforced & free text ending in
\texttt{[[A]]}, \texttt{[[B]]} or \texttt{[[C]]} \\
Confidence & reported by the judge & fixed at $0.85$ \\
\midrule
\multicolumn{3}{l}{\textit{Identical:} question set, stored responses, judge
model, dual-call reversal, tie-on-disagreement} \\
\bottomrule
\end{tabular}
\end{table}

The consequence is that any measured difference between the arms is attributable
to the bundle, not to the rubric alone. The reasoning scaffold runs in the
direction opposite to the observed effect --- the rubric arm is the one that
forces a four-step critique before a verdict --- so scaffolding is unlikely to
explain a generic-arm advantage, but the output format and the fixed confidence
are not controlled. Isolating the rubric would require running the MT-Bench
system prompt through the rubric arm's own user template and parser, which was
not done.

Under the replay design of Section~\ref{sec:exp-rqmap}, the two arms see the
same comparison set by construction. Measured overlap is $100\%$ on
mathematics, physics, psychology and philosophy, and $99.6\%$ on law and
engineering, where a handful of replays failed and were not retried.

\section{What the Judge Prompt Contains}\label{sec:judge-prompt-contents}

The mechanism behind this thesis's central RQ1 finding is visible in the
prompt, so what it does and does not contain is stated explicitly. In every
mode the judge receives the leaf's induced rubric as a system prompt, the
question stem, and the two stored responses; it never receives the MMLU-Pro
reference answer. The prompt-construction path asserts that the rendered query
carries no ground-truth answer block, and the stored correct answer is never
read on the judging path.

What varies is the option context.

\begin{description}
\item[\texttt{OPTIONS}] appends the full multiple-choice options block to the
stem.
\item[\texttt{RESOLVED}] withholds the options block and instead attaches each
model's own selected option to its trace as text, so a verdict can refer to an
answer rather than to a bare letter. This is the settled batch's mode.
Attachment succeeded on $93.2\%$ to $99.8\%$ of traces by domain; the shortfall
is traces whose upstream \texttt{pred} field is missing, which are passed
through verbatim.
\end{description}

The measured consequence is that ``answer-key-blind'' is a much weaker property
than ``answer-blind''. Under either mode a judge able to solve the item can
verify both responses against its own solution without ever being told the key,
and Chapter~\ref{ch:results} reports that this is what it does. Withholding the
options block narrows the channel but does not close it, because the responses
themselves state their answers.

Full source: \texttt{src/main/kotlin/taxonomy/prompts/JudgePrompts.kt}
and \texttt{src/main/kotlin/taxonomy/prompts/}\allowbreak
\texttt{GenericPairwiseJudgePrompt.kt}.

\section{A Worked Example}\label{sec:judge-example}

Everything above describes the judge; this section shows one. The prompts below
are reproduced mechanically from the frozen snapshot and the Kotlin sources by
\texttt{report/Figures/make\_judge\_example.py}, rather than transcribed, so
they cannot drift from what the runs actually sent.

% GENERATED by report/Figures/make_judge_example.py -- do not edit.
% Leaf n00000066 (Geometric Optics of Thin and Thick Lenses) of snapshot 20260727_042523_Headless_Run_Auto_ge.
\subsection{The Induced System Prompt}\label{sec:judge-example-system}

Leaf \texttt{n00000066}, \emph{Geometric Optics of Thin and Thick Lenses}. This half is generated once at snapshot-freeze time and is what the rubric arm sends as the judge's system prompt. It names the domain, its formalism and its failure modes; it says nothing about how to compare two responses.

\begin{quote}\small
You are an expert AI judge specializing in the geometric optics of thin and thick lenses, with deep mastery of sign conventions, lensmaker equations, image formation, magnification, refractive power, and system modeling. Your role is to evaluate responses to multiple-choice questions in this domain with precision, ensuring strict adherence to optical physics principles, correct application of formulas, and avoidance of common pitfalls. You prioritize the following: (1) Rigorous application of sign conventions (e.g., Cartesian, real-is-positive) for object/image distances and focal lengths, (2) Correct selection and derivation of lens/mirror equations (thin-lens formula, thick-lens matrix optics, or mirror equation), (3) Accurate propagation of image properties (position, size, orientation, real/virtual nature) through single or compound systems, (4) Proper handling of refractive power (diopters), media transitions (e.g., air to water), and system constraints (e.g., apertures, aberrations), (5) Distinction between linear and angular magnification, and (6) Physical plausibility of results (e.g., no negative focal lengths for converging lenses, no real images formed by diverging lenses). You identify and penalize failure modes such as inconsistent sign conventions, misapplied formulas, unit errors, unjustified approximations, and overlooked edge cases (e.g., object at focal point, infinite image distance). Your expertise extends to thick lenses (principal planes, equivalent focal lengths), multi-element systems (telescopes, corrective lenses), and aberrations (chromatic, spherical), but you strictly limit evaluation to geometric optics assumptions (paraxial rays, no wave effects).\par
\end{quote}
\subsection{The Induced Rubric}\label{sec:judge-example-rubric}

The rubric for the same leaf, reproduced in full. Its criteria are domain conditions --- sign conventions, unit consistency, physical plausibility --- and none of them is a comparison instruction, a tie rule, or an output format. That separation is the control described in Section~\ref{sec:judge-template-full}.

\begin{quote}\small
- **Sign Conventions and Formula Selection**: Response must explicitly state and consistently apply the correct sign convention (e.g., Cartesian, real-is-positive) for object distance, image distance, and focal length. Must derive or apply the appropriate lensmaker's equation, mirror equation, or thick-lens formula without mixing conventions or formulas.\par
- **Image Formation and Magnification**: Response must correctly classify images as real/virtual and erect/inverted, using consistent sign conventions. Must compute magnification (linear or angular) step-by-step for each optical element, propagating intermediate image positions and sizes accurately in compound systems.\par
- **Refractive Power and System Properties**: Response must correctly convert between refractive power (diopters) and focal length, ensuring unit consistency. Must adjust focal lengths and refractive powers for media transitions (e.g., air to water) and account for effective apertures or f-number changes.\par
- **System Modeling and Ray Tracing**: For thick lenses or multi-element systems, response must calculate equivalent focal lengths and principal plane positions or apply iterative thin-lens formulas. Must trace rays accurately through composite systems, including telescopes or corrective lenses.\par
- **Aberrations and Physical Constraints**: Response must account for dispersion (e.g., different refractive indices for wavelengths) or diffraction limits if relevant. Must ensure the effective f-number is adjusted for object distance and aperture size, especially when the lens is not focused at infinity.\par
- **Angular vs. Linear Magnification**: Response must distinguish between angular magnification (e.g., telescopes) and linear magnification (e.g., lenses/mirrors), applying the correct formula for the system.\par
- **Physical Plausibility**: Response must validate results for physical plausibility (e.g., positive focal lengths for converging lenses, virtual images not projected on screens). Must check edge cases (e.g., object at focal point, infinite image distance) to avoid invalid conclusions.\par
- **Unit Consistency and Dimensional Analysis**: Response must maintain consistent units (e.g., meters for distances, diopters for power) and perform dimensional analysis to catch errors (e.g., focal length in diopters vs. meters).\par
- **Approximations and Problem Interpretation**: Response must justify approximations (e.g., paraxial, small-angle) and explicitly state assumptions (e.g., coordinate system, reference points) to resolve ambiguities. Must not overlook deviations from exact ray tracing or wave optics effects.\par
- **Failure Mode Penalization**: Response must avoid common pitfalls: mixing sign conventions, misapplying formulas (e.g., thin-lens vs. thick-lens), ignoring media transitions, or misclassifying image types (real/virtual, erect/inverted). Errors in these areas are critical and must be penalized.\par
\end{quote}
\subsection{The Two User Messages}\label{sec:judge-example-user}

The rubric arm's user message, from \texttt{TaxonomyArenaService.buildJudgeUserPrompt}. The horizontal rule is a box-drawing character in the source and is shown here as hyphens.

\begin{quote}\small\ttfamily
{}[Question] \\
\$query \\
~ \\
{}[Model A's Response] \\
\$traceA \\
~ \\
{}[Model B's Response] \\
\$traceB \\
~ \\
---------------------------------------- \\
Your task is to determine superior reasoning quality based solely on the logical rigour and domain \\
soundness of its response. \\
~ \\
If both models are genuinely equivalent, output winner: \textquotedbl{}TIE\textquotedbl{} with confidence \textless{}= 0.5. \\
~ \\
Evaluate following the mandatory sequence: \\
~ \\
Step 1 -- Critique Model A (2–4 sentences referencing your rubric) \\
Step 2 -- Critique Model B (2–4 sentences referencing your rubric) \\
Step 3 -- Compare: what is the decisive difference in reasoning quality? \\
Step 4 -- Verdict \\
~ \\
Output ONLY the following JSON. No prose before or after. No markdown fences. \\
~ \\
\{ \\
  \textquotedbl{}critique\_a\textquotedbl{}: \textquotedbl{}\textless{}your critique of Model A\textgreater{}\textquotedbl{}, \\
  \textquotedbl{}critique\_b\textquotedbl{}: \textquotedbl{}\textless{}your critique of Model B\textgreater{}\textquotedbl{}, \\
  \textquotedbl{}comparison\textquotedbl{}: \textquotedbl{}\textless{}decisive difference in reasoning quality\textgreater{}\textquotedbl{}, \\
  \textquotedbl{}winner\textquotedbl{}: \textquotedbl{}Model A, Model B, or TIE\textquotedbl{}, \\
  \textquotedbl{}confidence\textquotedbl{}: \textless{}0.0 to 1.0 (default to 0.7-0.85 for clear wins, reserve 0.95-1.0 only for overwhelming evidence)\textgreater{} \\
\}
\end{quote}

The generic arm's user message, from \texttt{GenericPairwiseJudgePrompt.USER\_TEMPLATE}, is the MT-Bench pairwise template unchanged.

\begin{quote}\small\ttfamily
{}[User Question] \\
\{question\} \\
~ \\
{}[The Start of Assistant A's Answer] \\
\{answer\_a\} \\
{}[The End of Assistant A's Answer] \\
~ \\
{}[The Start of Assistant B's Answer] \\
\{answer\_b\} \\
{}[The End of Assistant B's Answer]
\end{quote}

Its system prompt is likewise the MT-Bench text verbatim.

\begin{quote}\small
Please act as an impartial judge and evaluate the quality of the responses provided by two AI assistants to the user question displayed below. You should choose the assistant that follows the user's instructions and answers the user's question better. Your evaluation should consider factors such as the helpfulness, relevance, accuracy, depth, creativity, and level of detail of their responses. Begin your evaluation by comparing the two responses and provide a short explanation. Avoid any position biases and ensure that the order in which the responses were presented does not influence your decision. Do not allow the length of the responses to influence your evaluation. Do not favor certain names of the assistants. Be as objective as possible. After providing your explanation, output your final verdict by strictly following this format: \textquotedbl{}[[A]]\textquotedbl{} if assistant A is better, \textquotedbl{}[[B]]\textquotedbl{} if assistant B is better, and \textquotedbl{}[[C]]\textquotedbl{} for a tie.\par
\end{quote}


\section{Rubric Reach: The Continuous Analysis}\label{sec:app-rubric-reach}

Demoted from Section~\ref{sec:res-offsubject}, which keeps the
headline numbers. Measured as cosine distance to the centroid of the
queries its rubric was induced from, a routed held-out question sits
at $0.511$ --- $38\%$ of the way from a random reserved question
($0.607$) toward construction-grade proximity ($0.353$), and only
$6$--$21\%$ of the way in five of eight domains. The premise of the
rubric-coverage worry is therefore real: rubrics are applied well
outside the material they were induced from. The consequence is not:
within cells, one standard deviation of that distance moves the
cell-scoped arm's tie rate by $+1.02$ percentage points --- and moves
the \emph{generic} arm, which carries no domain criteria at all, by
$+1.06$. The within-triple difference is $-0.04$ points
($95\%$ CI $[-1.39, +1.30]$, cell fixed effects, $n = 2{,}788$).

\begin{figure}[htbp]
\centering
\includegraphics[width=\textwidth]{rubric_reach}
\caption{Panel~(a): per domain, cosine distance from a routed
held-out question to its rubric's induction-pool centroid (dot),
between construction-grade proximity (left rule) and a random
held-out question (right rule). Panel~(b): change in tie rate and key
agreement per standard deviation of that distance, cell fixed
effects; the rubric-free arm moves identically, so the effect is not
rubric coverage.}
\label{fig:rubric-reach}
\end{figure}

\section{Dual-Call Position-Reversal Protocol}

Each logical comparison is realised as two judge calls with response
position reversed between calls, per the dual-call protocol of
Section~\ref{sec:arch-judge}. Disagreement between the two calls is
resolved as a tie, rather than by a coin flip or by discarding the
comparison. How that tie is then scored is a reporting convention and
not a modelling choice: every rank correlation appears under both
half-weighting and dropping. Making tie propensity a fitted parameter,
as \textcite{davidson1970ties} does, is the alternative this system
does not implement and is named as the principled repair
(Section~\ref{subsec:disc-construct-validity}).
